\section{Permutation-invariant energy parameterization}\label{app:perm-invariance}
FIME: Check notation and masking/padding
We implement the potential with a node/edge/global graph transformer backbone
(the GraphTransformer used in our codebase). Each layer is permutation equivariant:
for any relabeling $\sigma$, hidden node and edge features transform as
$(X,E)\mapsto (\sigma\!\cdot\!X,\sigma\!\cdot\!E)$. The global feature $h$ collects
extra graph-level features and is updated through permutation-invariant summaries,
so $h$ is invariant.

Here, $n$ denotes the maximum number of nodes observed in the training dataset; smaller graphs are padded and masked following standard conventions for graph diffusion models~\citep{vignac2022digress,qinmadeira2024defog}.

\begin{figure}[h]
\centering
\resizebox{\columnwidth}{!}{%
\begin{tikzpicture}[
  font=\scriptsize,
  node distance=1.0cm,
  module/.style={
    rectangle,
    rounded corners=5pt,
    minimum width=3.1cm,
    minimum height=0.95cm,
    align=center,
    draw=black!55,
    fill=black!4,
    thick
  },
  arrow/.style={-Latex, line width=1.1pt, draw=black!60}
]
  \node[module] (graph) {Graph input $x$\\(node + edge features)};
  \node[module, right=of graph] (gt) {Graph Transformer\\(equivariant)};
  \node[module, right=of gt] (pool) {Invariant readout\\$\sum_i X_i,\;\sum_{i\ne j} E_{ij}$ + \glsentryshort{mlp}};
  \node[module, right=of pool] (energy) {$V_\theta(x)\in\mathbb{R}$};

  \draw[arrow] (graph) -- (gt);
  \draw[arrow] (gt) -- (pool);
  \draw[arrow] (pool) -- (energy);
\end{tikzpicture}}
\caption{\textbf{Energy Backbone.} Permutation-equivariant backbone with invariant pooling to produce a scalar energy.}
\label{fig:energy-arch}
\end{figure}

The scalar energy is obtained by an invariant readout that first pools the
equivariant features and then applies an \gls{mlp}:
\begin{equation}
\begin{aligned}
\bar{X} &= \sum_{i\in\mathcal{V}} X^{(L)}_i,\qquad
\bar{E} = \sum_{i\ne j} E^{(L)}_{ij},\\
V_\theta(x,h) &= \mathrm{MLP}\big([\bar{X},\bar{E},h^{(L)}]\big).
\end{aligned}
\end{equation}
The sums are invariant to permutations, so the concatenated vector
$(\bar{X},\bar{E},h^{(L)})$ is invariant. Therefore any \gls{mlp} applied to it is
also invariant. This matches the sum-pooled energy head in the implementation
(with a masked off-diagonal sum).
